\section{Discussion, limitations and ethics}
\textbf{Recommendation for writers.} If the developer's pipeline is unknown, plant \emph{fictional facts phrased
differently each time}, about 16 copies of each across 100 or more facts (or 64 copies across 10 facts). That is
provable at $p<0.01$ with log-probability access after deduplication, normalization and paraphrasing. Watermarks
are far cheaper (1--4 copies) but rely on the pipeline skipping Unicode normalization, and verbatim canaries rely on
it skipping substring deduplication. With text-only access, none of our designs is reliable at realistic copy
counts. Making facts extractable by questioning, not just memorized as sentences, is the open problem (and the target
for optimized canaries).

\textbf{Pre-registered vs.\ stable rule.} Running 32 tests per model at $\alpha=0.01$ produces occasional isolated
hits. We saw three, all at 1 copy (Fig.~\ref{fig:dose}, hollow markers). The scorecard therefore counts a dose
only if every higher dose is also detected. This rule is exploratory, and Fig.~\ref{fig:dose} reports the
pre-registered one.

\textbf{Limitations.} One model size (160M) and one seed per pipeline; larger models memorize more, so our copy
counts are likely conservative \todo{unless tier 1 adds 410M}. Fluent traps are template-generated rather than
sampled from an LM at a controlled perplexity. The paraphraser rewrites only the canary paragraph, not the
surrounding document. Commercial pipelines also filter by quality and apply post-training, which we do not study
in the main results.

\textbf{Ethics.} We audit only models we trained. The protocol is designed so that a creator cannot falsely accuse
a model that never saw their text: the false-positive rate is controlled by the randomized split, which we verify
on the clean base model.
